\section[Conclusão]{Conclusão}

A AGES III foi uma experiência diferente das anteriores porque ocorreu em um momento em que eu já possuía mais maturidade técnica e experiência prática no mercado de software. Isso mudou a forma como enxerguei o projeto. Em vez de atuar apenas como alguém que estava aprendendo tecnologias novas, passei a perceber com mais clareza problemas de organização, comunicação, definição de escopo, gestão de pessoas e alinhamento entre squads.

No início do projeto, minha participação foi mais intensa. A proposta do VínculoHub Portal envolvia múltiplos perfis de usuário, autenticação, banco de dados, dashboards, documentos, editais, vínculos e uma arquitetura mais robusta. Esses elementos criaram oportunidades reais de contribuição, principalmente em decisões de produto, experiência do usuário, autenticação, organização de tarefas e apoio a colegas menos experientes. Nesse momento, senti que minha experiência anterior ajudava a dar mais segurança para o time.

Ao longo das sprints, entretanto, percebi uma redução gradual da minha presença e do meu engajamento no projeto. Essa mudança não esteve relacionada a uma dificuldade técnica específica, mas a uma combinação de fatores ligados ao meu interesse pelo produto e ao tipo de aprendizado que o projeto passou a oferecer. Nas primeiras etapas, havia desafios claros de estruturação, autenticação, organização de squads, definição de fluxos e apoio técnico. Com o avanço do semestre, muitas tarefas passaram a se concentrar em ajustes, refinamentos, revisão de pull requests, correções e sustentação de funcionalidades já encaminhadas.

Essa mudança me levou a uma autocrítica importante. Embora eu reconheça a relevância social da proposta do VínculoHub Portal, o domínio de negócio do projeto não despertava tanto interesse pessoal em mim. O business da plataforma, voltado à conexão entre organizações sociais e empresas, possui valor e impacto, mas não era um tema que me motivava profundamente. Como consequência, quando o projeto deixou de oferecer tantos aprendizados técnicos novos, ficou mais difícil sustentar o mesmo nível de presença e energia.

Essa percepção foi uma das lições mais relevantes da AGES III. Passei a entender que, em projetos de software, o interesse pelo domínio de negócio influencia muito mais a qualidade da participação do que eu imaginava. Quando o desenvolvedor se conecta com o problema que está sendo resolvido, tende a propor melhorias com mais naturalidade, tomar decisões melhores e manter o engajamento mesmo em fases menos interessantes do desenvolvimento. Quando essa conexão não existe, a atuação pode se tornar mais mecânica, ainda que exista capacidade técnica para contribuir.

Esse aprendizado se torna ainda mais importante no contexto atual de disseminação do uso de agentes e ferramentas de inteligência artificial no desenvolvimento de software. À medida que essas ferramentas passam a automatizar partes cada vez maiores da implementação, escrever código tende a deixar de ser o principal diferencial isolado do desenvolvedor. O entendimento do problema, do usuário, do contexto de negócio e das decisões de produto passa a ganhar mais importância. Saber o que construir, por que construir e como aquilo gera valor se torna tão importante quanto saber implementar.

Nesse sentido, a AGES III me mostrou que a escolha de um projeto não deve considerar apenas a stack técnica ou a oportunidade de praticar ferramentas específicas. É necessário avaliar também o quanto o problema de negócio desperta curiosidade e vontade de aprofundamento. Em um cenário em que agentes podem acelerar a produção de código, o envolvimento crítico com o business se torna um fator central para diferenciar uma entrega tecnicamente funcional de uma solução realmente relevante.

Do ponto de vista técnico, considero que o projeto contribuiu para consolidar conhecimentos em autenticação, integração entre frontend e backend, organização de banco de dados, revisão de código, testes, documentação e infraestrutura. Também tive contato com práticas importantes de qualidade, como uso de Docker, Flyway, Auth0, Swagger, testes automatizados e validações de sprint. Mesmo quando minha atuação se tornou menos intensa, ainda participei de revisões, planejamento de tarefas, apoio a colegas e implementação de funcionalidades específicas.

Do ponto de vista de soft skills, a maior lição foi sobre comunicação e gestão de pessoas. Em várias sprints, ficou claro que o maior risco não estava apenas na complexidade técnica, mas na dificuldade de alinhar expectativas entre AGES III, AGES IV e os demais membros dos squads. User stories pouco claras, responsabilidades ambíguas e canais de comunicação frágeis geraram ruídos que poderiam ter sido reduzidos com uma gestão mais presente e objetiva.

Concluo a AGES III com uma visão mais madura sobre desenvolvimento de software. A experiência não foi apenas sobre construir uma plataforma, mas sobre compreender melhor meus próprios critérios de motivação, meus limites de engajamento e a importância crescente do domínio de negócio na profissão. Tecnicamente, o projeto foi relevante; pessoalmente, ele foi ainda mais importante por me mostrar que escolher bem o problema a ser resolvido é parte essencial da construção de uma carreira sustentável em software.
